% Options for packages loaded elsewhere
\PassOptionsToPackage{unicode}{hyperref}
\PassOptionsToPackage{hyphens}{url}
\documentclass[
]{ltjsarticle}
\usepackage{xcolor}
\usepackage[margin=20mm]{geometry}
\usepackage{amsmath,amssymb}
\setcounter{secnumdepth}{-\maxdimen} % remove section numbering
\usepackage{iftex}
\ifPDFTeX
  \usepackage[T1]{fontenc}
  \usepackage[utf8]{inputenc}
  \usepackage{textcomp} % provide euro and other symbols
\else % if luatex or xetex
  \usepackage{unicode-math} % this also loads fontspec
  \defaultfontfeatures{Scale=MatchLowercase}
  \defaultfontfeatures[\rmfamily]{Ligatures=TeX,Scale=1}
\fi
\usepackage{lmodern}
\ifPDFTeX\else
  % xetex/luatex font selection
\fi
% Use upquote if available, for straight quotes in verbatim environments
\IfFileExists{upquote.sty}{\usepackage{upquote}}{}
\IfFileExists{microtype.sty}{% use microtype if available
  \usepackage[]{microtype}
  \UseMicrotypeSet[protrusion]{basicmath} % disable protrusion for tt fonts
}{}
\makeatletter
\@ifundefined{KOMAClassName}{% if non-KOMA class
  \IfFileExists{parskip.sty}{%
    \usepackage{parskip}
  }{% else
    \setlength{\parindent}{0pt}
    \setlength{\parskip}{6pt plus 2pt minus 1pt}}
}{% if KOMA class
  \KOMAoptions{parskip=half}}
\makeatother
\usepackage{longtable,booktabs,array}
\newcounter{none} % for unnumbered tables
\usepackage{calc} % for calculating minipage widths
% Correct order of tables after \paragraph or \subparagraph
\usepackage{etoolbox}
\makeatletter
\patchcmd\longtable{\par}{\if@noskipsec\mbox{}\fi\par}{}{}
\makeatother
% Allow footnotes in longtable head/foot
\IfFileExists{footnotehyper.sty}{\usepackage{footnotehyper}}{\usepackage{footnote}}
\makesavenoteenv{longtable}
\setlength{\emergencystretch}{3em} % prevent overfull lines
\providecommand{\tightlist}{%
  \setlength{\itemsep}{0pt}\setlength{\parskip}{0pt}}
\usepackage{bookmark}
\IfFileExists{xurl.sty}{\usepackage{xurl}}{} % add URL line breaks if available
\urlstyle{same}
\hypersetup{
  hidelinks,
  pdfcreator={LaTeX via pandoc}}

\author{}
\date{}

\begin{document}

\section{Paper 15: Projection onto xyztRQ --- Construction and Machine
Verification of the Map and Inverse Map from the Discrete Dual System to
Spacetime
Coordinates}\label{paper-15-projection-onto-xyztrq-construction-and-machine-verification-of-the-map-and-inverse-map-from-the-discrete-dual-system-to-spacetime-coordinates}

\textbf{Author}: Noriaki Kihara \textbf{Date}: June 12, 2026
\textbf{Version}: v1.1 (content and results unchanged from v1.0;
nomenclature unified to Paper 0.5's displacement-record definition ---
ledger/register/bookkeeping mapped to conserved quantity / invariant /
accounting) \textbf{DOI (this version)}: 10.5281/zenodo.20690060 /
\textbf{Concept DOI}: 10.5281/zenodo.20665688 \textbf{License}: CC BY
4.0 \textbf{Series}: Dual Geometry of Wavelength Space and Frequency
Space (sequel to Papers 1--14) \textbf{Note}: originally reviewed and
accepted as ``Paper 13''; renumbered to \textbf{Paper 15} upon insertion
of the new Papers 13 (position and time resolved) and 14 (time direction
and the amplitude stage) --- the only content change is the redirection
of citations. The numbering is independent of other series.

\begin{center}\rule{0.5\linewidth}{0.5pt}\end{center}

\subsection{Abstract}\label{abstract}

From the state data of the discrete dual system constructed in Papers
1--12 (three axioms --- reciprocal duality νλ=1, zero point ½,
asymmetric single bit --- plus the operating principles of configuration
reading and the record principle), we construct an \textbf{explicit
projection and inverse projection} onto the spacetime coordinates
\((x,y,z,t)\) and the invariant coordinates \((R,Q)\), and
machine-verify every stage. Main results: (1) \textbf{Spatial
coordinates} are given by a mixed-radix \([4;k,k,\dots]\) quarter-digit
expansion (\(k\) = the odd genealogy ratio); the digits are read from
the phase classes (Z₄) of the transcription lines of the record by
peeling decode. The only genealogy for which the readout closes within
the digital alphabet of the record is the \textbf{ternary one
(\(k{=}3\))} (information count \(2\,\mathrm{bit}\ge\log_2 k\)). (2)
\textbf{Time is a derived quantity}: \(\nu_t=\sqrt{s}\) (the norm
clock), \(t=\sum 1/\nu_t\) (event accumulation). It has no independent
inverse map --- not a defect, but the map-level realization of the
theorem (Paper 13, Theorem 6) that \(t\) is not a fundamental degree of
freedom. (3) \textbf{R and Q require no map}: \(R=\sqrt{s}\) and
\(Q=\varepsilon\) are B₄ invariants shared by all markings, and the
image of the projection lies on the constraint surface
\(\Sigma=\{R^2=s,\ Q=\varepsilon\}\) (codimension 2) --- the formal
statement of ``subjective 4D within six background axes.'' (4)
\textbf{Inverse map}: space is constructively invertible (rounding loss
below the censorship quantum; composite round trip 500/500), time admits
only event counting, R/Q are constraints. (5) \textbf{Compatibility with
conservation laws}: the value of \(s\) read from the record alone is
additive across decays (60/60 --- a check that could have failed), and
the multiplicative conservation of the charge-like \(Q\) is a
\textbf{theorem of parity arithmetic on the support} (Theorem 8),
enforced mechanically by transport: 39.2\% of raw cell combinations
violate ε, yet zero of the amplitude-carrying paths do. (6) \textbf{4D
lift}: product waves possess no pure-axis lines (sum/difference lines
only), and the simultaneous half-wave shift of an even number of active
axes is an identity of the state --- the readout decodes exactly to that
quotient (4-axis exhaustive: 32 classes, one-to-one), while parity-mixed
composites break the degeneracy completely (all 256 configurations
uniquely decoded). (7) \textbf{Kinematics}: uniform motion is a
mixed-radix odometer of the position digits; the amplitude spectrum,
wavelength content, R, and Q do not evolve in time (deviation
\(10^{-15}\)) --- only the phase evolves, linearly, and its rotation
rate is the velocity (exactly \(f\cdot v\)). \textbf{Velocity exists
nowhere in a single record; it exists only in the sequence of records.}
This paper does not claim a re-derivation of standard theory. No
physical identification is made.

\begin{center}\rule{0.5\linewidth}{0.5pt}\end{center}

\subsection{1. Introduction}\label{introduction}

\subsubsection{1.1 The question}\label{the-question}

Papers 1--12 constructed a discrete system from the reciprocal duality
\(\nu\lambda=1\), the zero point ½, and a single asymmetric bit, and
showed that position (Paper 13, Theorem 1 and Definition 2 / Theorems
2a--2b), time (Paper 13, Theorems 5--6), and the stage for amplitudes
(Paper 14, Theorems 1--3) are handled by the existing inventory alone.
What remained was the \textbf{map itself}:

\begin{quote}
Can the state data of the system be mapped to, and inverted from, the
standard physical coordinates \((x,y,z,t)\) and the invariant
coordinates \((R,Q)\) as an \textbf{actual operation}?
\end{quote}

There is a distance between ``it should be possible'' and ``the
formalization is exhibited and machine-checked.'' This paper closes that
distance: the projection is given as a five-stage operation, each stage
machine-verified (methods listed in Appendix A), and beyond ``values
round-trip'' we demonstrate on the output side that the map
\textbf{commutes with the conservation laws} --- that R is a derived
quantity governed by an energy-like conserved quantity and that t is
derived from R.

\subsubsection{1.2 Position within the
series}\label{position-within-the-series}

This paper is the \textbf{coordinate-map installment} of the
correspondence-dictionary arc (Paper 12, §8); the measure line is
outside its scope and remains in active development --- the arc is not
yet complete. The transport principle --- ``theorems of standard theory
transport without re-derivation when their premises have verified
counterparts in this system'' --- becomes operational only when the
coordinate map connecting the two sides has been constructed. The
materials are the state structure of Papers 5--7, the internal time of
Paper 8, the odd-nesting theorem of Paper 9, and the readout theory of
Paper 12; no new axiom is added.

\subsubsection{1.3 What is not claimed}\label{what-is-not-claimed}

We do not claim (i) a re-derivation of standard theory, (ii) a
construction of common time for multi-fragment systems
(synchronization), or (iii) a derivation of dynamics (why uniform motion
occurs). (ii) belongs to the gauge-fixing mechanism of the record
interface (a separate line); (iii) to the effective interpolation of the
wave sector (Paper 12, §4.2). No physical identification is made.

\subsection{2. Preliminaries (minimal
recap)}\label{preliminaries-minimal-recap}

State data: occupied cells \(k\in\mathbb{Z}^4\) (configuration reading:
the state \emph{is} the set of occupied cells = the wave itself), branch
(cos/sin) and sign per level, the genealogy chain (odd-ratio nesting
\(k_\ell\), Paper 9), and the conserved quantities
\(s=\sum_i(|k_i|+\tfrac12)^2\), \(\varepsilon=(-1)^{\Sigma|k_i|}\). The
record: the hologram \(I=\Psi^2\) --- its physical content is the
\textbf{line coefficients (relational data)}; the position-space picture
is their Pontryagin-dual chart (space = the character group of the
conserved-charge lattice; Paper 13, Theorems 1--2). Gauge data of the
chart: the marking \(u\) (choice of the time-read axis; 16 choices
forming a single B₄ orbit (Paper 14, Theorem 2)) and the
\textbf{per-axis orientations} \(Z_2^{\,4}\) (choice of positive branch
direction --- the position-reading counterpart of the orientation fixing
of Paper 14, §3). These are held fixed below.

\textbf{Emphasis}: a particle does not ``remember'' its state. Position
= the phase offsets of its wave; R = the norm of its wavelength content;
Q = the parity of its content --- all are aspects that \emph{constitute}
the state. The only memory in the system is the record; every readout in
this paper takes the record alone as input.

\subsection{3. Definition of the
projection}\label{definition-of-the-projection}

Fix a marking \(u\). The three transverse axes provide \((x,y,z)\).

\textbf{D1 (spatial coordinates)}: for each axis \(j\neq u\), the
mixed-radix quarter-digit expansion

\[
x_j=\frac{1}{4}\Bigl[d_j^{(0)}+\sum_{\ell\ge1}c_j^{(\ell)}\prod_{m\le\ell}k_m^{-1}\Bigr]\bmod 1,
\qquad d^{(0)}\in Z_4,\ c^{(\ell)}\in Z_{k_\ell},\ k_\ell\ \text{odd}
\]

where \(d^{(0)}\) is the level-0 quarter digit (cos=0, sin=1, −cos=2,
−sin=3) and \(c^{(\ell)}\) the position of the level-\(\ell\) child comb
within the parent quarter gap.

\textbf{D2 (time coordinate)}: \(t=\sum_{\text{record events}}1/\nu_t\)
with \(\nu_t=\sqrt{s}\) --- identical in construction to the internal
time of Paper 8, §3.2 (numerically cross-checked).

\textbf{D3 (R, Q)}: \(R=\sqrt{s}=\nu_t\), \(Q=\varepsilon\). No map
required (identity reads); complete B₄ invariants shared by all
markings.

\textbf{D4 (inverse map)}: below we specialize the genealogy to the
canonical \(k_\ell=3\) (Theorem 4). For \(x_j\): greedy rounding to
depth \(L\),

\[
d^{(0)}=\lfloor4x\rfloor\bmod4,\qquad
c^{(\ell)}=\Bigl\lfloor 3^{\ell}\bigl(4x-d^{(0)}-\sum_{1\le m<\ell}c^{(m)}3^{-m}\bigr)\Bigr\rfloor\bmod 3
\]

with the error guarantee
\(|x-\hat x_L|\le\frac14(\prod k_m)^{-1}\frac{k_L}{k_L-1}\) (below the
censorship quantum = no loss of physical fact). In implementation,
floating-point boundary points may yield a digit equal to \(k\); we clip
with \(\min(c,k{-}1)\) (a measure-zero boundary convention that does not
affect the bound). For \(t\): no independent inverse (only the event
count \(n=\lfloor t\sqrt s\rfloor\)). For \(R,Q\): constraints (no
freedom).

\subsection{4. Basic theorems}\label{basic-theorems}

\textbf{Theorem 1 (quarter-digit consistency)}: the translation
\(T_{1/4}\) shifts the digit \(d\to d+1\pmod4\) (exhaustive, 4/4).

\textbf{Theorem 2 (completeness and rounding bound)}: the chain of
position lattices
\(P_\ell=(4k_1\cdots k_\ell)^{-1}\mathbb{Z}/\mathbb{Z}\) is nested for
any odd genealogy (\([P_\ell:P_{\ell-1}]=k_\ell\)); D1 is a complete
mixed-radix expansion on \(T\) (20,000 random points within the bound).
An important distinction: \textbf{the hierarchical scale of position is
the wavelength genealogy ratio \(k\), not the container ratio
\(1/(2R')\)} (the latter makes the digit chain lacunary --- rejected by
the machine). Oddness originates not from density but from wave
consistency (the odd-nesting theorem).

\textbf{Theorem 3 (constraint surface)}: the image of the projection is
the codimension-2 surface

\[
\Sigma=\bigl\{(x,y,z,t,R,Q): R^2=s(k),\ Q=\varepsilon(k)\bigr\}.
\]

R and Q carry no inverse-map freedom --- the formalization of the
``invisible directions'' (Paper 12, Suppl. 39/40). Virtual displacements
(±1 borrowing) are off-surface moves permitted only between records
(§7).

\subsection{5. The readout operation (five stages of the forward
map)}\label{the-readout-operation-five-stages-of-the-forward-map}

For axis \(j\) and frequency \(f=3^\ell\) (henceforth \(k=3\)), the
transcription line:

\begin{enumerate}
\def\labelenumi{\arabic{enumi}.}
\tightlist
\item
  \textbf{Line extraction}:
  \(Z^{(\ell)}=2\langle I,\cos(2\pi 3^\ell x_j)\rangle+2i\,\langle I,\sin(2\pi 3^\ell x_j)\rangle\)
\item
  \textbf{Peeling}:
  \(\tilde Z^{(\ell)}=Z^{(\ell)}-\sum_{m<\ell}(\text{known transcription contributions of decoded shallower levels})\).
  \textbf{Even-drop theorem}: comb×comb interference falls entirely on
  even frequencies (odd×odd = even; leakage \(2\times10^{-15}\)), so odd
  readout lines consist of transcription terms only and the peeled
  amounts are exactly known.
\item
  \textbf{Quarter class}:
  \(r^{(\ell)}=\frac{2}{\pi}\arg\tilde Z^{(\ell)}\bmod 4\in Z_4\) (all
  phases land exactly on the quarter lattice; quantization error
  \(3.4\times10^{-15}\))
\item
  \textbf{Digit decode (recursion)}: \(c^{(0)}=r^{(0)}\),
  \(c^{(\ell)}=(r^{(\ell)}-3r^{(\ell-1)})\bmod4\)
\item
  \textbf{Assembly}: substitute into D1.
\end{enumerate}

Verification: two levels 12/12; three levels 36/36; full combs (all odd
harmonics, cutoff 81) three levels 36/36 --- all uniquely decoded.

\textbf{Theorem 4 (digital preference for the ternary genealogy)}: a
quarter channel carries \(Z_4=2\) bits per level. Unique decoding of a
digit \(c\in Z_k\) requires \(k\le4\); among odd \(k\ge3\) only \(k=3\)
qualifies (\(k{=}1\) is a degenerate genealogy and excluded; \(k{=}5\)
exhibits the degeneracy \(c{=}4\equiv c{=}0\), machine-demonstrated
4/20). \textbf{The only genealogy whose positions can be read out
entirely within the record's digital alphabet is the ternary cascade}
--- coinciding, by an independent principle, with the genealogy selected
by minimality in Paper 8 (the Z₂ prohibition of binary).

\subsection{6. The 4D lift}\label{the-4d-lift}

A product wave \(\Phi=\prod_j\varphi(x_j-a_j)\) has no pure-axis lines
(its spectrum consists of sum/difference lines only); digits are
obtained by solving the quarter classes of sum/difference lines mod 4.

\textbf{Theorem 5 (diagonal quotient)}: the simultaneous half-wave shift
of an even number of active axes (\(d_j\to d_j+2\)) is an
\textbf{identity of the wave} (2-axis 16/16; 4-axis: state classes =
quotient by the even-sign-flip group \(Z_2^3\), 256/8 = 32). The readout
decodes exactly to this quotient (32 signatures, one-to-one) --- a
redundancy of coordinates, not an ambiguity of states (configuration
reading: same wave = same state).

\textbf{Theorem 6 (full decoding via parity mixing)}: composites
containing cells with an odd number of active axes break the degeneracy.
The condition is a \textbf{covering of every nontrivial even flip by an
odd-parity cell intersecting it oddly}; constructions satisfying it
decode uniquely in full enumeration --- 2-axis 16/16 and \textbf{4-axis
256/256}.

\subsection{7. Compatibility with conservation laws --- R and t are
derived
quantities}\label{compatibility-with-conservation-laws-r-and-t-are-derived-quantities}

Of the six coordinates xyztRQ, only \(xyz\) plus the event count are
free. The checks in this section fall into three distinct strata (made
explicit at referee point \#2):

\textbf{(i) Accounting consistency (true by construction)}: the decay
path table enforces \(s_{\rm kept}+s_c+s_d=s_{\rm parent}\) at records
and the ±1 borrowing cancellation at virtual stages by definition. Its
holding over all 118,944/118,944 contributing paths at s=9/11/13 (eight
representative parents per s, covering all three orbit types) is a
\textbf{consistency check that the path enumeration and the amplitude
machinery implement the same accounting} --- it has no refutation power.

\textbf{(ii) A theorem (which could have failed, and is proven)}:

\begin{quote}
\textbf{Theorem 8 (charge-like parity selection rule; the general-lemma
form is canonical as Paper 14, Appendix A-3 --- this theorem is its
selection-rule expression)}: if a transport coefficient is nonzero,
\(\varepsilon\) is multiplicatively conserved. Proof: \(C_c\neq0\)
requires the target \(v_p=n_a+n_b\) (with \(n\) on the coefficient
support, \(|n_{a,j}|=|v_{a,j}|\)), and integer-sum parity gives
\(|v_{p,j}|\equiv|v_{a,j}|+|v_{b,j}|\pmod 2\) per axis; summing over
axes, \(\varepsilon(v_p)=\varepsilon(v_a)\varepsilon(v_b)\). Applying
this to both stages,
\(\varepsilon_{\rm parent}=\prod\varepsilon_{\rm daughters}\). ∎
\end{quote}

Machine confirmation: zero violations over all 118,944 contributing
paths. Control: among raw cell combinations (s=13, 11.4 million)
\textbf{39.2\% violate ε} --- violating candidates abound on the
support, and the readout dictionary annihilates their amplitudes
exactly. Charge-like conservation appears neither as a postulate nor as
an empirical selection effect, but as a \textbf{theorem of parity
arithmetic on the support}.

\textbf{(iii) A check that could have failed (via the record)}: the
occupied cell set is recovered from the hologram alone (exact
configuration reading, 60/60) → additivity of \(s_{\rm read}\) (60/60) →
the chain ``configuration → \(R=\sqrt{s_{\rm read}}\) → \(\nu_t=R\) →
\(t=\Sigma1/\nu_t\)'' holds at every stage \textbf{on the output side}.
This is independent of construction (the decoding could have failed; the
additivity could have broken).

\begin{quote}
The projection is not a picture of coordinates but a \textbf{faithful
display of the conserved quantities}. The conservation laws and the
projection are not merely compatible --- they are two faces of the same
dictionary.
\end{quote}

\subsection{8. The inverse map and the composite round
trip}\label{the-inverse-map-and-the-composite-round-trip}

Continuous \(x\) → (D4 rounding) → digits → \textbf{actual construction
of the state} (full combs, three levels) → (§5 forward map) → digits,
\(\hat x\): over 500 random points, \textbf{complete digit agreement
500/500}, maximum error 0.0277 ≤ bound 0.0417.

\(t\): no independent inverse exists --- \(\nu_t\) is derived from the
spatial configuration (the realization of the theorem of Paper 13
(Theorem 6); D2 is its implementation). This is consistent with the
principle that ``t merely appears to have been selected by observation''
(the all-axes-symmetric principle). \(R,Q\): six-tuples violating the
constraints lie outside the image (= the invisible directions).

\subsection{9. Condition matrix and
robustness}\label{condition-matrix-and-robustness}

{\def\LTcaptype{none} % do not increment counter
\begin{longtable}[]{@{}
  >{\raggedright\arraybackslash}p{(\linewidth - 4\tabcolsep) * \real{0.3333}}
  >{\raggedright\arraybackslash}p{(\linewidth - 4\tabcolsep) * \real{0.3333}}
  >{\raggedright\arraybackslash}p{(\linewidth - 4\tabcolsep) * \real{0.3333}}@{}}
\toprule\noalign{}
\begin{minipage}[b]{\linewidth}\raggedright
Scenario
\end{minipage} & \begin{minipage}[b]{\linewidth}\raggedright
Configurations
\end{minipage} & \begin{minipage}[b]{\linewidth}\raggedright
Result
\end{minipage} \\
\midrule\noalign{}
\endhead
\bottomrule\noalign{}
\endlastfoot
Three heterogeneous fragments (fundamentals 1/3/5; peeling cascade) & 64
& \textbf{64/64} \\
Three fragments at continuous positions (no digit alignment; direct
phase read) & 200 & \textbf{200/200} (error \(2.5\times10^{-16}\)) \\
Harmonic amplitude jitter ±20\% + additive noise σ=0.5 & 192 &
\textbf{192/192} \\
Three same-species fragments (worst case: shared lines) & \textbf{all
220} & without half-wave pair \textbf{160/160} unique / with
\textbf{60/60} state-identical (Theorem 7) \\
Three 4D composites (configuration reading) & 60 & \textbf{60/60} \\
\end{longtable}
}

\textbf{Theorem 7 (half-wave pair = state identity)}: when a
same-species pair sits at half-wavelength offset, the sum of the square
combs vanishes identically (maximum \(1.5\times10^{-14}\)) --- i.e.,
\textbf{a configuration containing such a pair is the same wave, hence
the same state, as the configuration without it} (like Theorem 5, an
element of the kernel of the configuration→wave map). The readout
returns the state correctly; the ``degeneracy'' is mere label
redundancy. This separates the present theorem from the open problem of
§12-2 (\(k\ge5\) deep digits = \textbf{different states} that cannot be
read = censorship): the former is state identity, the latter a limit of
readability.

\subsection{10. Kinematics: the map of uniform
motion}\label{kinematics-the-map-of-uniform-motion}

This section verifies uniform motion by \textbf{two complementary
constructions} (they are not a single object --- referee point \#10):

\textbf{O1 (digit odometer: readout of hierarchical states)}: each tick
(\(\Delta t=1/\sqrt s\)) advances the deepest quarter quantum by one,
and the \textbf{hierarchical state, re-assembled digit-aligned at every
tick}, is read. The hierarchical decode returns digit = counter value
(base {[}4;3,3{]}) at all 36 ticks; \(\hat x_n=n/36\) exactly. The world
line is an exact straight line through lattice points,
\(v=\Delta a\cdot\sqrt s\).

\textbf{O2 (rigid drift: spectrum of a single comb)}: a single comb is
rigidly translated continuously and the time evolution of its line
spectrum is measured:

\begin{itemize}
\tightlist
\item
  \textbf{The amplitude spectrum, wavelength content, R, and Q do not
  evolve} (maximum deviation \(1.9\times10^{-15}\) over 36 ticks)
\item
  Only the phase evolves: exactly linear rotation at rate \(f\cdot v\)
  (residual \(\sim10^{-15}\); ratio \(f{=}3/f{=}1=3.000000\))
\end{itemize}

The two are distinct verifications: a rigidly drifting hierarchical
composite leaves the quarter lattice at intermediate ticks, and in
digit-aligned staircase motion the phases of shallow lines advance
stepwise (exact linearity holds only for the deepest line). The ``rigid
drift'' of Figure 4 refers to O2.

\begin{quote}
\textbf{Velocity appears nowhere in a single snapshot} --- both
\textbar spectrum\textbar{} and support are identical to the static
state; velocity exists only in the sequence of records. On the
conserved-quantity side: \textbar spectrum\textbar{} = the conserved
quantity of free motion (momentum-like), phase = position, phase
rotation rate = velocity. Readable velocities lie on a rational lattice;
continuous velocity is its completion (the same architecture as for
position).
\end{quote}

\subsection{11. Figures}\label{figures}

\begin{itemize}
\tightlist
\item
  \textbf{Figure 1} (\texttt{paper15\_fig1\_readout.png}): the readout
  in both spaces --- transcription lines in frequency space (odd lines
  only), superposed combs in position space, fragments on the constraint
  surface in xyztRQ space, condition-matrix results
\item
  \textbf{Figure 2} (\texttt{paper15\_fig2\_xyztRQ.png}): xyztRQ space
  --- the \((x,y,R)\) 3D arrangement (positions machine-decoded), world
  lines (tick spacing \(1/R\), branching at decay), the \((R,Q)\)
  invariant plane, the constraint surface \(R^2=s\) with virtual ±1
\item
  \textbf{Figure 3} (\texttt{paper15\_fig3\_conservation.png}):
  conservation --- spectra before/after the decay \(9\to(5,1,3)\), the
  \(R^2\) conserved quantity (virtual +1 borrowing), the ε control
  (39.2\% vs 0), verification summary
\item
  \textbf{Figure 4} (\texttt{paper15\_fig4\_uniform\_motion.png}):
  uniform motion --- the x--t world line, the spacetime record (rigid
  drift, O2), spectral invariance, linear phase rotation
\end{itemize}

All figures are exact computations (no schematic drawings).

\subsection{12. Limits and outlook}\label{limits-and-outlook}

\begin{enumerate}
\def\labelenumi{\arabic{enumi}.}
\tightlist
\item
  \textbf{Common time for multiple fragments}: D2 is the proper clock of
  a single fragment. Synchronization = the (1,3) aggregation (gauge
  fixing by the record interface) is not yet constructed --- the
  system's most important remaining item.
\item
  \textbf{Analog digits for \(k\ge5\)}: deep digits of non-ternary
  genealogies require sub-quarter amplitude-ratio reads, outside the
  digital alphabet. Whether ``an unreadable digit exists as position''
  touches the principle of configuration reading (an open problem; the
  canonical genealogy is ternary, so the main line is unaffected).
\item
  \textbf{Dynamics}: the motion of this paper is kinematics (the map and
  readout of a schedule); deriving ``no force → uniform motion'' belongs
  to the wave-sector interpolation. Acceleration, relative motion, and
  collisions are not constructed.
\item
  General-\(s\) algebraic proofs of single-sign support, torsor
  residues, etc. share the series' proof queue (s=9/11/13 are exhaustive
  facts).
\item
  No physical identification is made. Standard theory is exactly
  correct; this paper constructs the map portion of ``a different
  mapping of the same results.''
\end{enumerate}

\subsection{Appendix A: Verification summary (all
reproducible)}\label{appendix-a-verification-summary-all-reproducible}

{\def\LTcaptype{none} % do not increment counter
\begin{longtable}[]{@{}
  >{\raggedright\arraybackslash}p{(\linewidth - 6\tabcolsep) * \real{0.2500}}
  >{\raggedright\arraybackslash}p{(\linewidth - 6\tabcolsep) * \real{0.2500}}
  >{\raggedright\arraybackslash}p{(\linewidth - 6\tabcolsep) * \real{0.2500}}
  >{\raggedright\arraybackslash}p{(\linewidth - 6\tabcolsep) * \real{0.2500}}@{}}
\toprule\noalign{}
\begin{minipage}[b]{\linewidth}\raggedright
Claim
\end{minipage} & \begin{minipage}[b]{\linewidth}\raggedright
Verification
\end{minipage} & \begin{minipage}[b]{\linewidth}\raggedright
Method
\end{minipage} & \begin{minipage}[b]{\linewidth}\raggedright
Script
\end{minipage} \\
\midrule\noalign{}
\endhead
\bottomrule\noalign{}
\endlastfoot
Quarter-digit shift consistency & 4/4 & exhaustive &
supplement62\_projection\_formalization.py \\
Mixed-radix completeness and bound & 20,000 pts, k=3,5 & random &
same \\
Constraint-surface invariance & B₄ samples & sampled & same \\
t = internal time of Paper 8 & numerical match & cross-check & same \\
Readout (fundamental combs) & 12/12, 36/36 & exhaustive &
supplement63\_digit\_readout\_protocol.py \\
k=5 degeneracy (preference theorem) & 4/20 & exhaustive & same \\
Full combs, peeling & 36/36, err 3.4e-15 & exhaustive &
supplement63\_full\_comb\_collision\_test.py \\
Even drop & leakage 2e-15 & exhaustive (cutoff 81) & same \\
2-axis diagonal quotient, parity mixing & 16/16, 8 classes, 144, 16/16 &
exhaustive & supplement64\_4d\_product\_wave\_test.py \\
\textbf{4-axis full check} & 32 classes one-to-one, \textbf{256/256} &
exhaustive & paper15\_appendix\_4axis\_check.py \\
Composite round trip (inverse) & 500/500 & random &
supplement62\_inverse\_roundtrip\_test.py \\
Accounting consistency (R² additivity, ±1) & 118,944/118,944 &
exhaustive (true by construction, §7-i) &
supplement65\_conservation\_tests.py \\
Theorem 8 (ε selection rule), machine confirmation & 118,944/118,944;
control 39.2\%→0 & exhaustive + theorem & supplement65\_controls.py \\
Configuration reading, s\_read additivity & 60/60 & random (final
states) & supplement65\_c3\_exact\_decoder.py \\
Condition matrix B/C/D & 64/64, 200/200, 192/192 & exh./random/random &
supplement66\_condition\_matrix.py \\
Same-species classification & 160/160 + 60/60 & \textbf{exhaustive
(220)} & same + this paper \\
Uniform motion O1/O2 & 36/36, 1.9e-15, exact f·v & exhaustive (36 ticks)
& supplement67\_uniform\_motion.py \\
\end{longtable}
}

\begin{center}\rule{0.5\linewidth}{0.5pt}\end{center}

\textbf{Acknowledgments / history}: the verification of this paper was
carried out under the two-party independent verification protocol of
Claude Code (local machine verification) and claude.ai (independent
re-computation and review). Source material: Supplements 62--67 (June
12, 2026); foundational theorems now in Papers 13--14. v0.2: 14 referee
points incorporated (review with independent re-implementation
reproducing all computational claims); v1.0: 3 residual points of the
second-round review (accept).

No physical identification is made.

\end{document}
